\documentclass[11pt]{article}
\newcommand{\ListaTitulo}{Lista 01 --- Princípios do Delineamento}
% Preâmbulo compartilhado — MATD48, Listas de Exercícios 2026
% Usado por Lista{NN}.tex e Gabarito{NN}.tex via % Preâmbulo compartilhado — MATD48, Listas de Exercícios 2026
% Usado por Lista{NN}.tex e Gabarito{NN}.tex via % Preâmbulo compartilhado — MATD48, Listas de Exercícios 2026
% Usado por Lista{NN}.tex e Gabarito{NN}.tex via \input{preamble}

\usepackage[utf8]{inputenc}
\usepackage[T1]{fontenc}
\usepackage[brazil]{babel}
\usepackage{fullpage}
\usepackage{amsmath,amssymb}
\usepackage{enumitem}
\usepackage{xcolor}
\usepackage{fancyhdr}
\usepackage{titlesec}
\usepackage{listings}
\usepackage{hyperref}
\usepackage{booktabs}
\usepackage{graphicx}

\definecolor{matd48azul}{HTML}{0B4F6C}
\definecolor{matd48verde}{HTML}{2E8B57}

\titleformat{\section}{\normalfont\Large\bfseries\color{matd48azul}}{\thesection}{1em}{}
\titleformat{\subsection}{\normalfont\large\bfseries\color{matd48azul}}{\thesubsection}{1em}{}

\providecommand{\ListaTitulo}{Lista de Exercícios}

\setlength{\headheight}{14pt}
\pagestyle{fancy}
\fancyhf{}
\lhead{\small MATD48 --- Planejamento de Experimentos A}
\rhead{\small \ListaTitulo}
\cfoot{\thepage}
\renewcommand{\headrulewidth}{0.4pt}

\lstset{
  basicstyle=\ttfamily\small,
  breaklines=true,
  frame=single,
  columns=fullflexible,
  backgroundcolor=\color{gray!8},
  commentstyle=\color{matd48verde},
  keywordstyle=\color{matd48azul}\bfseries,
  language=R,
  extendedchars=true,
  literate=%
    {á}{{\'a}}1 {à}{{\`a}}1 {â}{{\^a}}1 {ã}{{\~a}}1
    {é}{{\'e}}1 {ê}{{\^e}}1 {í}{{\'i}}1 {ó}{{\'o}}1
    {ô}{{\^o}}1 {õ}{{\~o}}1 {ú}{{\'u}}1 {ç}{{\c c}}1
    {Á}{{\'A}}1 {À}{{\`A}}1 {Â}{{\^A}}1 {Ã}{{\~A}}1
    {É}{{\'E}}1 {Ê}{{\^E}}1 {Í}{{\'I}}1 {Ó}{{\'O}}1
    {Ô}{{\^O}}1 {Õ}{{\~O}}1 {Ú}{{\'U}}1 {Ç}{{\c C}}1
}

\hypersetup{colorlinks=true, linkcolor=matd48azul, urlcolor=matd48azul, citecolor=matd48azul}

\newcommand{\cabecalho}[2]{%
  \begin{center}
    {\Large\bfseries\color{matd48azul} #1}\\[0.3em]
    {\large #2}\\[0.3em]
    \rule{\linewidth}{0.6pt}
  \end{center}
  \vspace{0.5em}
}

\newenvironment{questao}[1]{\vspace{1em}\noindent\textbf{Questão #1.}\hspace{0.4em}}{\par}
\newenvironment{solucao}{\vspace{0.3em}\noindent\textit{Solução.}\hspace{0.4em}\color{black!85}}{\par\vspace{0.5em}}


\usepackage[utf8]{inputenc}
\usepackage[T1]{fontenc}
\usepackage[brazil]{babel}
\usepackage{fullpage}
\usepackage{amsmath,amssymb}
\usepackage{enumitem}
\usepackage{xcolor}
\usepackage{fancyhdr}
\usepackage{titlesec}
\usepackage{listings}
\usepackage{hyperref}
\usepackage{booktabs}
\usepackage{graphicx}

\definecolor{matd48azul}{HTML}{0B4F6C}
\definecolor{matd48verde}{HTML}{2E8B57}

\titleformat{\section}{\normalfont\Large\bfseries\color{matd48azul}}{\thesection}{1em}{}
\titleformat{\subsection}{\normalfont\large\bfseries\color{matd48azul}}{\thesubsection}{1em}{}

\providecommand{\ListaTitulo}{Lista de Exercícios}

\setlength{\headheight}{14pt}
\pagestyle{fancy}
\fancyhf{}
\lhead{\small MATD48 --- Planejamento de Experimentos A}
\rhead{\small \ListaTitulo}
\cfoot{\thepage}
\renewcommand{\headrulewidth}{0.4pt}

\lstset{
  basicstyle=\ttfamily\small,
  breaklines=true,
  frame=single,
  columns=fullflexible,
  backgroundcolor=\color{gray!8},
  commentstyle=\color{matd48verde},
  keywordstyle=\color{matd48azul}\bfseries,
  language=R,
  extendedchars=true,
  literate=%
    {á}{{\'a}}1 {à}{{\`a}}1 {â}{{\^a}}1 {ã}{{\~a}}1
    {é}{{\'e}}1 {ê}{{\^e}}1 {í}{{\'i}}1 {ó}{{\'o}}1
    {ô}{{\^o}}1 {õ}{{\~o}}1 {ú}{{\'u}}1 {ç}{{\c c}}1
    {Á}{{\'A}}1 {À}{{\`A}}1 {Â}{{\^A}}1 {Ã}{{\~A}}1
    {É}{{\'E}}1 {Ê}{{\^E}}1 {Í}{{\'I}}1 {Ó}{{\'O}}1
    {Ô}{{\^O}}1 {Õ}{{\~O}}1 {Ú}{{\'U}}1 {Ç}{{\c C}}1
}

\hypersetup{colorlinks=true, linkcolor=matd48azul, urlcolor=matd48azul, citecolor=matd48azul}

\newcommand{\cabecalho}[2]{%
  \begin{center}
    {\Large\bfseries\color{matd48azul} #1}\\[0.3em]
    {\large #2}\\[0.3em]
    \rule{\linewidth}{0.6pt}
  \end{center}
  \vspace{0.5em}
}

\newenvironment{questao}[1]{\vspace{1em}\noindent\textbf{Questão #1.}\hspace{0.4em}}{\par}
\newenvironment{solucao}{\vspace{0.3em}\noindent\textit{Solução.}\hspace{0.4em}\color{black!85}}{\par\vspace{0.5em}}


\usepackage[utf8]{inputenc}
\usepackage[T1]{fontenc}
\usepackage[brazil]{babel}
\usepackage{fullpage}
\usepackage{amsmath,amssymb}
\usepackage{enumitem}
\usepackage{xcolor}
\usepackage{fancyhdr}
\usepackage{titlesec}
\usepackage{listings}
\usepackage{hyperref}
\usepackage{booktabs}
\usepackage{graphicx}

\definecolor{matd48azul}{HTML}{0B4F6C}
\definecolor{matd48verde}{HTML}{2E8B57}

\titleformat{\section}{\normalfont\Large\bfseries\color{matd48azul}}{\thesection}{1em}{}
\titleformat{\subsection}{\normalfont\large\bfseries\color{matd48azul}}{\thesubsection}{1em}{}

\providecommand{\ListaTitulo}{Lista de Exercícios}

\setlength{\headheight}{14pt}
\pagestyle{fancy}
\fancyhf{}
\lhead{\small MATD48 --- Planejamento de Experimentos A}
\rhead{\small \ListaTitulo}
\cfoot{\thepage}
\renewcommand{\headrulewidth}{0.4pt}

\lstset{
  basicstyle=\ttfamily\small,
  breaklines=true,
  frame=single,
  columns=fullflexible,
  backgroundcolor=\color{gray!8},
  commentstyle=\color{matd48verde},
  keywordstyle=\color{matd48azul}\bfseries,
  language=R,
  extendedchars=true,
  literate=%
    {á}{{\'a}}1 {à}{{\`a}}1 {â}{{\^a}}1 {ã}{{\~a}}1
    {é}{{\'e}}1 {ê}{{\^e}}1 {í}{{\'i}}1 {ó}{{\'o}}1
    {ô}{{\^o}}1 {õ}{{\~o}}1 {ú}{{\'u}}1 {ç}{{\c c}}1
    {Á}{{\'A}}1 {À}{{\`A}}1 {Â}{{\^A}}1 {Ã}{{\~A}}1
    {É}{{\'E}}1 {Ê}{{\^E}}1 {Í}{{\'I}}1 {Ó}{{\'O}}1
    {Ô}{{\^O}}1 {Õ}{{\~O}}1 {Ú}{{\'U}}1 {Ç}{{\c C}}1
}

\hypersetup{colorlinks=true, linkcolor=matd48azul, urlcolor=matd48azul, citecolor=matd48azul}

\newcommand{\cabecalho}[2]{%
  \begin{center}
    {\Large\bfseries\color{matd48azul} #1}\\[0.3em]
    {\large #2}\\[0.3em]
    \rule{\linewidth}{0.6pt}
  \end{center}
  \vspace{0.5em}
}

\newenvironment{questao}[1]{\vspace{1em}\noindent\textbf{Questão #1.}\hspace{0.4em}}{\par}
\newenvironment{solucao}{\vspace{0.3em}\noindent\textit{Solução.}\hspace{0.4em}\color{black!85}}{\par\vspace{0.5em}}


\begin{document}

\cabecalho{Lista de Exercícios 01}{Princípios do delineamento de experimentos (Capítulo 1 / Aula 01)}

\noindent \textit{Entrega: até a próxima aula. Justifique todas as respostas — nestas listas, a
justificativa vale mais que a resposta final. O gabarito completo está em \texttt{Gabarito01.pdf}.}

\begin{questao}{1} \textbf{(Psicologia --- identificação de unidades)}

Um pesquisador estuda o efeito de três programas de meditação guiada (10, 20 e 30 minutos
diários) sobre o nível de ansiedade autorreportado. Ele recruta 24 participantes, sorteia 8 para
cada programa e, ao final de quatro semanas, aplica o mesmo questionário de ansiedade \emph{três
vezes} a cada participante, em dias diferentes da mesma semana.

\begin{enumerate}[label=(\alph*)]
  \item Identifique a unidade experimental e a unidade amostral deste estudo.
  \item Quantas réplicas legítimas existem para comparar os três programas? Justifique.
  \item Se o pesquisador rodar uma ANOVA usando as $24 \times 3 = 72$ respostas diretamente,
    que problema estatístico ele comete?
\end{enumerate}
\end{questao}

\begin{questao}{2} \textbf{(Agricultura --- fator, nível, tratamento)}

Uma equipe agronômica testa o efeito de duas variáveis sobre a produtividade do feijão: a
variedade da semente (V1, V2, V3) e a densidade de plantio (baixa, alta). Todas as combinações
de variedade e densidade são usadas.

\begin{enumerate}[label=(\alph*)]
  \item Quantos fatores existem neste estudo? Liste os níveis de cada um.
  \item Quantos tratamentos distintos resultam do cruzamento dos fatores?
  \item Esse é um desenho de tratamentos fatorial ou de um único fator? Explique a diferença.
\end{enumerate}
\end{questao}

\begin{questao}{3} \textbf{(Ciência de dados --- os três pilares)}

Uma equipe de produto quer saber se um novo algoritmo de recomendação aumenta o tempo de sessão
dos usuários em um aplicativo de música. Eles implementam o algoritmo novo apenas para usuários
que atualizaram o aplicativo na última semana, e comparam o tempo de sessão desse grupo com o de
usuários que não atualizaram.

\begin{enumerate}[label=(\alph*)]
  \item Este desenho usa aleatorização? Justifique.
  \item Aponte pelo menos duas variáveis de confusão plausíveis que podem explicar uma diferença
    observada, mesmo que o algoritmo não tenha efeito real.
  \item Reformule o desenho do experimento para que ele permita uma conclusão causal válida,
    especificando fator, níveis, unidade experimental e o mecanismo de aleatorização.
\end{enumerate}
\end{questao}

\begin{questao}{4} \textbf{(Fontes de variação)}

Para cada cenário abaixo, classifique a fonte de variação descrita como (i) efeito do
tratamento, (ii) variação sistemática entre unidades experimentais, ou (iii) erro experimental.

\begin{enumerate}[label=(\alph*)]
  \item A diferença de peso ganho entre dois grupos de frangos que receberam rações diferentes.
  \item A diferença de peso inicial entre os frangos, antes de qualquer ração ser oferecida.
  \item A variação no peso final entre frangos do \emph{mesmo} grupo, alimentados com a mesma
    ração, nascidos da mesma ninhada.
\end{enumerate}
\end{questao}

\begin{questao}{5} \textbf{(Controle local / blocagem --- dissertativa curta)}

Um professor quer comparar dois métodos de ensino (expositivo vs. baseado em projetos) e sabe,
de antemão, que suas quatro turmas do semestre têm desempenhos historicamente muito diferentes
entre si.

\begin{enumerate}[label=(\alph*)]
  \item Proponha um desenho que use blocagem por turma. Descreva como a aleatorização
    aconteceria \emph{dentro} de cada bloco.
  \item Explique, em suas palavras, por que blocar por turma tende a tornar a comparação entre
    métodos mais precisa do que sortear os métodos ignorando a turma.
  \item A blocagem substitui a necessidade de aleatorização? Justifique.
\end{enumerate}
\end{questao}

\begin{questao}{6} \textbf{(Dedução --- resultados potenciais)}

Considere $N$ unidades experimentais com resultados potenciais fixos $Y_1(0),\dots,Y_N(0)$ e
$Y_1(1),\dots,Y_N(1)$ (SUTVA vale). Um mecanismo de aleatorização sorteia, com probabilidade
uniforme entre todos os subconjuntos de tamanho $n_1$, quais $n_1$ unidades recebem $Z_i=1$; as
$n_0=N-n_1$ restantes recebem $Z_i=0$. Define-se
$\widehat{\text{ATE}} = \bar Y_1^{\,obs} - \bar Y_0^{\,obs}$, com
$\bar Y_1^{\,obs} = \tfrac{1}{n_1}\sum_{i:Z_i=1} Y_i(1)$ e analogamente para $\bar Y_0^{\,obs}$.

\begin{enumerate}[label=(\alph*)]
  \item Mostre que $\bar Y_1^{\,obs}$ é a média de uma amostra aleatória simples sem reposição, de
    tamanho $n_1$, retirada do conjunto fixo $\{Y_1(1),\dots,Y_N(1)\}$.
  \item Usando o fato de que a média de uma AAS sem reposição é um estimador não-viesado da média
    populacional, mostre que $E\big[\widehat{\text{ATE}}\big] = \text{ATE} = \bar Y(1) - \bar Y(0)$.
  \item Explique, em uma frase, por que essa demonstração \emph{não} usa nenhuma suposição sobre a
    distribuição de $Y_i(0)$ ou $Y_i(1)$ (por exemplo, não exige normalidade).
\end{enumerate}
\end{questao}

\begin{questao}{7} \textbf{(Método científico --- falseabilidade e validade externa)}

Uma startup de e-commerce testa um novo algoritmo de recomendação apenas com os usuários que se
inscreveram voluntariamente em seu "programa de testadores beta" (um grupo pequeno, engajado, que
recebe recompensas por participar). Dentro desse grupo, os usuários são corretamente sorteados
entre "recomendação antiga" e "recomendação nova" — a aleatorização é impecável.

\begin{enumerate}[label=(\alph*)]
  \item No vocabulário do Capítulo 1, essa aleatorização impecável garante \emph{validade
    interna}, \emph{validade externa}, ou as duas? Justifique.
  \item Reformule a alegação de resultado da empresa ("a recomendação nova aumenta as vendas em
    12\%") como uma hipótese \emph{falseável}, no sentido de Popper — isto é, descreva que
    resultado observável, se encontrado, contradiria a alegação.
  \item A diretoria quer generalizar a conclusão do teste beta para \emph{todos} os usuários do
    site, não só os testadores voluntários. Que suposição adicional (não garantida só pela
    aleatorização dentro do grupo beta) essa generalização exige?
\end{enumerate}
\end{questao}

\end{document}
